%document
\documentclass[10pt]{beamer}
%theme
\usetheme{metropolis}
% packages
\usepackage{color}
\usepackage{listings}
\usepackage[ngerman]{babel}
\usepackage[utf8x]{inputenc}
\usepackage{multicol}


% color definitions
\definecolor{mygreen}{rgb}{0,0.6,0}
\definecolor{mygray}{rgb}{0.5,0.5,0.5}
\definecolor{mymauve}{rgb}{0.58,0,0.82}

\lstset{
    backgroundcolor=\color{white},
    % choose the background color;
    % you must add \usepackage{color} or \usepackage{xcolor}
    basicstyle=\footnotesize\ttfamily,
    % the size of the fonts that are used for the code
    breakatwhitespace=false,
    % sets if automatic breaks should only happen at whitespace
    breaklines=true,                 % sets automatic line breaking
    captionpos=b,                    % sets the caption-position to bottom
    commentstyle=\color{mygreen},    % comment style
    % deletekeywords={...},
    % if you want to delete keywords from the given language
    extendedchars=true,
    % lets you use non-ASCII characters;
    % for 8-bits encodings only, does not work with UTF-8
    frame=single,                    % adds a frame around the code
    keepspaces=true,
    % keeps spaces in text,
    % useful for keeping indentation of code
    % (possibly needs columns=flexible)
    keywordstyle=\color{blue},       % keyword style
    % morekeywords={*,...},
    % if you want to add more keywords to the set
    numbers=left,
    % where to put the line-numbers; possible values are (none, left, right)
    numbersep=5pt,
    % how far the line-numbers are from the code
    numberstyle=\tiny\color{mygray},
    % the style that is used for the line-numbers
    rulecolor=\color{black},
    % if not set, the frame-color may be changed on line-breaks
    % within not-black text (e.g. comments (green here))
    stepnumber=1,
    % the step between two line-numbers.
    % If it's 1, each line will be numbered
    stringstyle=\color{mymauve},     % string literal style
    tabsize=4,                       % sets default tabsize to 4 spaces
    % show the filename of files included with \lstinputlisting;
    % also try caption instead of title
    language = Java,
	showspaces = false,
	showtabs = false,
	showstringspaces = false,
	escapechar = ,
}

\def\ContinueLineNumber{\lstset{firstnumber=last}}
\def\StartLineAt#1{\lstset{firstnumber=#1}}
\let\numberLineAt\StartLineAt



\newcommand{\codeline}[1]{
	\alert{\texttt{#1}}
}

% This Document contains the information about this course.

% Authors of the slides
\author{Fritz Windisch, Lukas Tassanyi}

% Name of the Course
\institute{Java-Course}

% Fancy Logo
\titlegraphic{\hfill\includegraphics[height=1.25cm]{../templates/fsr_logo_cropped}}


\title{Java}
\subtitle{Controll Statements and OOP}
\date{\today}

\def\code#1{\texttt{#1}}

\begin{document}

\begin{frame}
    \titlepage
\end{frame}
\begin{frame}{Overview}
    \setbeamertemplate{section in toc}[sections numbered]
    \tableofcontents
\end{frame}

\section{Recalling last session}
\begin{frame}{Recalling last session}
    Datatypes
    \begin{itemize}
        \item int, long
        \item float, double
        \item String
    \end{itemize}
    Hello World example
\end{frame}

\section{Controll Statements}
\begin{frame}{Controll Statements}
    
    \begin{itemize}
        \item if, else, else if
        \item for
        \item while
    \end{itemize}
        
\end{frame}

\subsection{\code{if} and \code{else}}
\begin{frame}[fragile]{\code{if} and \code{else}}
\begin{lstlisting}
if(condition) {
    // do something if condition is true
} else if(another condition){
    // do if "else if" condition is true 
} else {
    // otherwise do this
}\end{lstlisting}
\end{frame}

\begin{frame}[fragile]{\code{if} and \code{else} example}
\begin{lstlisting}
public class IteExample {

    public static void main(String[] args) {
        int myNumber = 5;
        
        if(myNumber >= 10) {
            System.out.println("Your number is 10 or greater.");
        } else if(myNumber == 3) {
            System.out.println("Your number is a 3.");
        } else {
            System.out.println("Your number is smaller than 10 but not a 3.");
        }
    }
    
}\end{lstlisting}
\end{frame}

\begin{frame}{Conditions 1/2}

Conditions can be either \code{true} of \code{false}.
\\[0.3in]

Comparisons:
\begin{itemize}
    \item \texttt{a == b} ... equal
    \item \texttt{a != b} ... not equal
    \item \texttt{a > b} ... greater than
    \item \texttt{a >= b} ... greater or equal than
\end{itemize}
\end{frame}

\begin{frame}{Conditions 2/2}

Conditions can be concatenated to more complex conditions:
\\[0.3in]

Logical Operations:
\begin{itemize}
    \item \texttt{x \&\& y} ... x and y
    \item \texttt{x || y} ... x or y
    \item \texttt{!x} ... not x
\end{itemize}
\end{frame}

\subsection{\code{while}}
\begin{frame}[fragile]{\code{while}}
\begin{lstlisting}
while(condition) {
    // do code while condition is true
}\end{lstlisting}
\end{frame}

\begin{frame}[fragile]{\code{while} example}
\begin{lstlisting}
public class WhileExample {

    public static void main(String[] args) {
        int i = 0;
        while(i <= 10) {
            System.out.println(i);
            i++;
        }
    }
    
}\end{lstlisting}
\end{frame}

\subsection{\code{for}}
\begin{frame}[fragile]{\code{for}}
\begin{lstlisting}
int i = 0;
while(i <= 10) {
    // do something ... 
    i++;
}

for(int i = 0; i <= 10; i++) {
    // do something ... 
}\end{lstlisting}
\end{frame}

\begin{frame}[fragile]{\code{for} example}
\begin{lstlisting}
public class ForExample {

    public static void main(String[] args) {
        for(int i = 0; i <= 10; i++) {
            System.out.print("na ");
        }
        System.out.println("BATMAN!");
    }
    
}\end{lstlisting}
\end{frame}


\section{OOP in Java}

\begin{frame}{}
    \begin{center}
        {\huge Object Oriented Programming}
    \end{center}
\end{frame}

\section{Conclusion}
\subsection{An Example}

\begin{frame}{An Example}
    You want to program an enrollment system, for a programming course. \\
    \vspace{1em}
    Your classes are:\\
    \begin{description}
        \item[student] who wants to attend the course
        \item[lesson] which is a part of the course
        \item[tutor] the guy with the bandshirt
        \item[room] where your lessons take place
        \item[\dots]
    \end{description}
%     The more you think about it, the more complex this program becomes.
%     Focus on the relevant things.
%    Think about how the objects can be in relation, this will be discussed later
%    Show prepared classes in Java
\end{frame}

\subsection{General information}

\begin{frame}[fragile]{Class Student}
\begin{lstlisting}
public class Student {

    // Attributes
    private String name; 
    private int matriculationNumber; 
    
    
    // Methods
    public void setName(String name) {
        this.name = name;
    }
    
    public int getMatriculationNumber() {
        return matriculationNumber;
    }
    
}\end{lstlisting}

% What is visible here:
% Attributes store the state of the object
% Methods implement the behaviour of the object

\end{frame}

\begin{frame}[fragile]{Creation}
    We learned how to declare and assign a primitive datatype.
    
    \begin{lstlisting}
        int a; // declare a
        a = 273; // assign 273 to a \end{lstlisting} 
    
    The creation of an object works similar.
    
    \begin{lstlisting}
        Student example = new Student(); 
        // create an instance of Student \end{lstlisting}
    The \textbf{object} derived from a \textbf{class} is also called \textbf{instance}.
    The variable is called the \textbf{reference}.
\end{frame}

\subsection{Methods}
\begin{frame}[fragile]{Calling a Method}
    \begin{lstlisting}
    public class Student {
    
        private String name;
    
        public String getName() {
            return name;
        }
        
        public void setName(String newName) {
            name = newName;
        }
       
    } \end{lstlisting}
    The class \emph{Student} has two methods: \emph{void getName()} and \emph{void setName()}.
\end{frame}

\begin{frame}[fragile]{Calling a Method}
\begin{lstlisting}
public class Main {

    public static void main(String[] args) {
        Student example = new Student(); // creation
        example.setName("Jane"); // method call
        String name = example.getName(); 
        System.out.println(name); // Prints "Jane"
    }
    
}\end{lstlisting}
    You can call a method of an object after its creation with \textbf{reference.methodName();}.
\end{frame}

\begin{frame}[fragile]{Calling a Method}
\begin{lstlisting}
    public class Student {
    
        private String name;
    
        public void setName(String newName) {
            name = newName;
            printName();   // Call own method
            this.printName(); // Or this way
        }
        
        public void printName() {
            System.out.println(name);
        }
        
    } \end{lstlisting}
    You can call a method of the own object by simply writing \textbf{methodName();} or \textbf{this.methodName();}
\end{frame}

\begin{frame}[fragile]{Methods with Arguments}

\begin{lstlisting}
public class Calc {

    public void add(int summand1, int summand2) {
        System.out.println(summand1 + summand2);
    }
        
    public static void main(String[] args) {
        int summandA = 1;
        int summandB = 2;
        Calc calculator = new Calc();
        System.out.print("1 + 2 = ");
        calculator.add(summandA, summandB); 
        // prints: 3
    }
    
} \end{lstlisting}
\end{frame}

\subsection{Return Value}
\begin{frame}[fragile]{Methods with Return Value}
    A method without a return value is indicated by \textbf{void}:
    \begin{lstlisting}
    public void add(int summand1, int summand2) {
        System.out.println(summand1 + summand2);
    }\end{lstlisting}
    A method with an \textbf{int} as return value:
    \begin{lstlisting}
    public int add(int summand1, int summand2) {
        return summand1 + summand2;
    } \end{lstlisting}
    % TODO explain return statement
\end{frame}

\begin{frame}[fragile]{Calling Methods with a return value}
    \begin{lstlisting}
    public class Calc {
    
        public int add(int summand1, int summand2) {
            return summand1 + summand2;
        }
        
        public static void main(String[] args) {
            Calc calculator = new Calc();
            int sum = calculator.add(3, 8);
            System.out.print("3 + 8 = " + sum); 
            // prints: 3 + 8 = 11
        }
        
    } \end{lstlisting}
\end{frame}

\subsection{Constructor}

\begin{frame}[fragile]{Constructors}
    \begin{lstlisting}
    public class Calc {
    
        private int summand1;
        private int summand2;
    
        public Calc() {
            summand1 = 0;
            summand2 = 0;
        }
        
    } \end{lstlisting}
    A constructor gets called upon creation of the object
\end{frame}

\begin{frame}[fragile]{Constructors with Arguments}
    \begin{lstlisting}
    public class Calc {
    
        private int summand1;
        private int summand2;
    
        public Calc(int x, int y) {
            summand1 = x;
            summand2 = y;
        }
        
    } \end{lstlisting}
    \begin{lstlisting}
    [...]
    Calc myCalc = new Calc(7, 9); \end{lstlisting}
    
    A constructor can have arguments as well!
\end{frame}

\begin{frame}[fragile]{Exercise and questions}
    \begin{itemize}
    \item You can now begin solving the exercises 2a-c from the course materials yourself or with a partner. \\
        \url{https://friwi.me/java-lessons/}
    \item If you require any assistance, feel free to ask us at any time
    \end{itemize}
\end{frame}

\begin{frame}[fragile]{Thank you for your attention}
    \begin{itemize}
    \item Course materials: \\
        \url{https://friwi.me/java-lessons/}
    \end{itemize}
\end{frame}

\end{document}
